% Источник решений: «семинар 2.pdf», страницы 1--4.
\section{Неклассические модели}

\noindent Условия:
\href{run:conditions/sem_02.pdf}{семинар 2 (PDF)}.

\subsection{Задача 1}

\begin{solution}
    \textbf{Модель: геометрическая.}
    $x$ --- опоздание семинариста, $y$ --- опоздание старосты.
    Выбор равномерный и независимый.
    \par\noindent
    \begin{minipage}[c]{0.62\linewidth}
        \[
            \omega=(x,y),\qquad \Omega=[0;15]\times[0;10].
        \]
        \[
            \mathcal{A}=\{(x,y)\in\Omega:x-10\leq y\leq x+5\}.
        \]
        \[
            S_\Omega=150,\qquad S_1=S_2=\frac{5\cdot5}{2}.
        \]
        \[
            \Pr(\mathcal{A})=\frac{150-S_1-S_2}{150}
            =\frac{125}{150}=\boxed{\frac56}.
        \]
    \end{minipage}\hfill
    \begin{minipage}[c]{0.36\linewidth}
        \centering
        \begin{tikzpicture}[x=0.25cm,y=0.25cm,>=Stealth]
            \fill[definition_frame!25]
                (0,0)--(10,0)--(15,5)--(15,10)--(5,10)--(0,5)--cycle;
            \draw (0,0) rectangle (15,10);
            \draw[thick] (0,5)--(5,10) (10,0)--(15,5);
            \draw[->] (0,0)--(17,0) node[right] {$x$};
            \draw[->] (0,0)--(0,12) node[above] {$y$};
            \node[below left] at (0,0) {$0$};
            \foreach \x in {5,10,15}
                \draw (\x,0.15)--(\x,-0.15) node[below] {$\x$};
            \foreach \y in {5,10}
                \draw (0.15,\y)--(-0.15,\y) node[left] {$\y$};
            \node at (1.4,8.5) {$S_1$};
            \node at (13.6,1.5) {$S_2$};
            \node at (7.5,5) {$\mathcal{A}$};
        \end{tikzpicture}
    \end{minipage}
\end{solution}

\subsection{Задача 2}

\begin{solution}
    \textbf{Модель: геометрическая.}
    $\omega=(\alpha,\beta)$, $\Omega=[0;1]^2$;
    $\eta$ --- число различных действительных корней.
    \[
        \begin{gathered}
            f_{\alpha,\beta}(x)=\frac{x^3}{3}-\alpha^2x+\beta,
            \qquad f'_{\alpha,\beta}(x)=x^2-\alpha^2;\\[4pt]
            f_{\alpha,\beta}(-\alpha)=\frac{2\alpha^3}{3}+\beta\geq0,
            \qquad f_{\alpha,\beta}(\alpha)=-\frac{2\alpha^3}{3}+\beta;\\[4pt]
            \eta=3
            \iff \alpha>0,\quad 0\leq\beta<\frac{2\alpha^3}{3}.
        \end{gathered}
    \]
    \par\begingroup\centering
        \begin{tikzpicture}[baseline=(current bounding box.center),
            x=1.25cm,y=1.25cm,>=Stealth]
            \draw[->] (-1.7,0)--(1.8,0) node[right] {$x$};
            \draw[->] (0,-0.65)--(0,0.85) node[above] {$f$};
            \draw[thick,domain=-1.6:1.6,samples=90]
                plot (\x,{\x^3/3-0.64*\x+0.1});
            \draw[dashed] (-0.8,0)--(-0.8,{2*0.8^3/3+0.1});
            \draw[dashed] (0.8,0)--(0.8,{-2*0.8^3/3+0.1});
            \node[below] at (-0.8,0) {$-\alpha$};
            \node[above] at (0.8,0) {$\alpha$};
        \end{tikzpicture}
        \hspace{1.2cm}
        \begin{tikzpicture}[baseline=(current bounding box.center),
            x=2.4cm,y=2.4cm,>=Stealth]
            \fill[definition_frame!25] (0,0)
                plot[domain=0:1,samples=60] (\x,{2*\x^3/3})
                --(1,1)--(0,1)--cycle;
            \draw (0,0) rectangle (1,1);
            \draw[thick,domain=0:1,samples=60]
                plot (\x,{2*\x^3/3});
            \draw[->] (0,0)--(1.15,0) node[right] {$\alpha$};
            \draw[->] (0,0)--(0,1.15) node[above] {$\beta$};
            \node[below left] at (0,0) {$0$};
            \node[below] at (1,0) {$1$};
            \node[left] at (0,1) {$1$};
            \node[right] at (1,0.6667) {$\frac23$};
            \node at (0.4,0.7) {$\eta=1$};
            \node[font=\small] at (0.77,0.12) {$\eta=3$};
        \end{tikzpicture}
    \par\endgroup
    Граница $\beta=2\alpha^3/3$ имеет нулевую площадь:
    $\Pr(\eta=2)=0$.
    \[
        \Pr(\eta=3)=\int_0^1\frac{2\alpha^3}{3}\diff\alpha
        =\boxed{\frac16},\qquad
        \Pr(\eta=1)=1-\frac16=\boxed{\frac56}.
    \]
\end{solution}

\clearpage
\subsection{Задача 3}

\begin{solution}
    \textbf{Три геометрические модели.}
    $\rho$ --- расстояние от центра до хорды, $\ell$ --- её длина.
    \[
        \ell=2\sqrt{R^2-\rho^2}>R\sqrt3
        \iff \rho<\frac R2.
    \]

    \textbf{1) Равномерный выбор конца хорды на окружности.}
    Один конец фиксирован; $s$ --- длина дуги до второго конца.
    \[
        \Omega=[0;2\pi R),\qquad
        \mathcal{A}=\left(\frac{2\pi R}{3};\frac{4\pi R}{3}\right),
        \qquad \Pr(\mathcal{A})=\frac{2\pi R/3}{2\pi R}
        =\boxed{\frac13}.
    \]

    \textbf{2) Равномерный выбор середины хорды в круге.}
    \[
        \Omega=\{(u,v):u^2+v^2\leq R^2\},\qquad
        \mathcal{A}=\{(u,v):u^2+v^2<R^2/4\}.
    \]
    \[
        \Pr(\mathcal{A})=\frac{\pi(R/2)^2}{\pi R^2}
        =\boxed{\frac14}.
    \]
    % При выборе центра направление диаметра задаётся произвольно:
    % вероятность выбора этой точки равна нулю.

    \textbf{3) Равномерный независимый выбор расстояния и угла.}
    \[
        \omega=(\rho,\varphi),\qquad
        \Omega=[0;R]\times[0;2\pi),\qquad
        \mathcal{A}=[0;R/2)\times[0;2\pi).
    \]
    \[
        \Pr(\mathcal{A})=\frac{(R/2)\cdot2\pi}{R\cdot2\pi}
        =\boxed{\frac12}.
    \]

    \par\begingroup\centering
        \begin{tikzpicture}[scale=1.25,>=Stealth]
            \draw (0,0) circle (1);
            \draw (1,0)--(-0.5,0.866)--(-0.5,-0.866)--cycle;
            \draw[definition_frame,very thick]
                (120:1) arc[start angle=120,end angle=240,radius=1];
            \fill (1,0) circle (0.025) node[right] {$A$};
            \fill (0,0) circle (0.02);
            \node at (0,-1.35) {1)};
        \end{tikzpicture}
        \hspace{1cm}
        \begin{tikzpicture}[scale=1.25]
            \fill[definition_frame!25] (0,0) circle (0.5);
            \draw (0,0) circle (1);
            \draw[thick] (0,0) circle (0.5);
            \draw (0,0)--(0.5,0) node[midway,above] {$\frac R2$};
            \fill (0,0) circle (0.02);
            \node at (0,-1.35) {2)};
        \end{tikzpicture}
        \hspace{1cm}
        \begin{tikzpicture}[scale=1.25,>=Stealth]
            \draw (0,0) circle (1);
            \draw[dashed] (0,0)--(1.1,0);
            \draw[->] (0.42,0) arc[start angle=0,end angle=35,radius=0.42];
            \node[below] at (0.48,-0.01) {$\varphi$};
            \begin{scope}[rotate=35]
                \draw[very thick,definition_frame] (0,0)--(0.65,0)
                    node[midway,above] {$\rho$};
                \draw[thick] (0.65,-0.7599)--(0.65,0.7599);
                \draw (0.54,0)--(0.54,0.11)--(0.65,0.11);
            \end{scope}
            \fill (0,0) circle (0.02);
            \node at (0,-1.35) {3)};
        \end{tikzpicture}
    \par\endgroup
\end{solution}

\subsection{Задача 4}

\begin{solution}
    \textbf{Модель: схема Бернулли $\mathrm{Bern}(2n,p)$.}
    \[
        \omega=(i_1,\ldots,i_{2n}),\qquad
        \Omega=\{0,1\}^{2n},\qquad
        \Pr(\{\omega\})=
        p^{\Sum{j=1}{2n}i_j}(1-p)^{2n-\Sum{j=1}{2n}i_j}.
    \]
    \[
        \mathcal{A}=
        \left\{\omega:i_2=i_4=\cdots=i_{2n}=1,\
        \Sum{j=1}{2n}i_j=n+m\right\}.
    \]
    В $n$ нечётных позициях должны стоять $m$ единиц:
    \[
        |\mathcal{A}|=C_n^m,\qquad
        \Pr(\mathcal{A})=
        \boxed{C_n^m p^{n+m}(1-p)^{n-m}},
        \qquad 0\leq m\leq n.
    \]
\end{solution}
